% sec-04: the weekly cadence.  Figure 14 (diagrams-type clustered topology),
% Table 23.
\section{The Weekly Cadence}

The block just completed was designed from nothing, which cost more time in
design than in execution. A week that produces five days of correspondence and no
recurring rhythm has to be reinvented, so the block closes by fixing one.

\begin{apptable}
\begin{tabularx}{\textwidth}{>{\raggedright\arraybackslash}p{2.2cm} >{\raggedright\arraybackslash}p{3.0cm} >{\raggedright\arraybackslash}p{1.6cm} >{\raggedright\arraybackslash}X}
\headrow \thead{Weekday} & \thead{Theme} & \thead{Budget} & \thead{The artifact that must exist by the end of it}\\
\midrule
Monday & Federal & 90 min & Sent letters, and any registration confirmed current \\
Tuesday & Capital & 75 min & A capital decision recorded, or a written reason none was taken \\
Wednesday & Site and partner & 90 min & Sent letters, and any meeting confirmed with an agenda attached \\
Thursday & Preparation & 120 min & The data room updated, and next week's drafts written and held \\
Friday & Execution and record & 75 min & Orders entered against their checks, follow-ups sent, the record table filed \\
\midrule
\textbf{Total} & Five standing functions & \textbf{7.5 h} & One person, sustainable alongside the technical work \\
\end{tabularx}
\end{apptable}
\vspace{-0.60cm}
\tabcap{Table 23. The weekly cadence by weekday, with a time budget and the one\\
artifact that must exist by the end of each day for the week to have held.}

\begin{appfloat}[!tb]
\begin{appfig}[diagrams-type / clustered topology with vector glyphs]
% Five clusters on a 2.85cm pitch, one glyph tile each.  \dgnode places the
% tile, its pictogram and its label as three related nodes, so a braced fit over
% the tile and its label encloses both.  The carry edge runs 1.4cm below the row.
\dgnode{dgtile}{0}{0}{\glyphbank}{Federal: re-contacts, program questions}{mo}
\dgnode{dgtile}{2.85}{0}{\glyphchart}{Capital: instruments, counsel, treasury}{tu}
\dgnode{dgtile}{5.70}{0}{\glyphteam}{Site and partner: institutions, foundations}{we}
\dgnodew{dgtilem}{8.55}{0}{\glyphdoc}{Preparation: data room, next week's drafts}{th}
\dgnode{dgtile}{11.40}{0}{\glyphclock}{Execution: release, follow-ups, record}{fr}
% \pnote is a macro, not a TikZ style; the budget labels below are plain nodes.
\foreach \x/\b in {0/{90 min}, 2.85/{75 min}, 5.70/{90 min}, 8.55/{120 min}, 11.40/{75 min}}{
  \node[font=\tiny\sffamily\bfseries,text=fundaccent,anchor=north] at (\x,-1.66) {\b};}
\foreach \x/\t/\n in {0/{Monday}/c1, 2.85/{Tuesday}/c2, 5.70/{Wednesday}/c3,
                      8.55/{Thursday}/c4, 11.40/{Friday}/c5}{
  \node[dgctitle,anchor=south] (\n) at (\x,0.62) {\t};}
\node[dgcluster,fit={(mo)(mol)(c1)},inner sep=7pt] {};
\node[dgcluster,fit={(tu)(tul)(c2)},inner sep=7pt] {};
\node[dgcluster,fit={(we)(wel)(c3)},inner sep=7pt] {};
\node[dgcluster2,fit={(th)(thl)(c4)},inner sep=7pt] {};
\node[dgcluster,fit={(fr)(frl)(c5)},inner sep=7pt] {};
\draw[dgedgeb] (mo) -- (tu);
\draw[dgedgeb] (tu) -- (we);
\draw[dgedgeb] (we) -- (th);
\draw[dgedgeb] (th) -- (fr);
\draw[dgedged] (11.40,-2.30) -- (11.40,-2.90) -- (0,-2.90) -- (0,-2.30);
\node[font=\tiny\itshape,text=fundgray,anchor=north] at (5.70,-2.98)
  {Carry rule: unfinished work moves to the same weekday next week, never to the next day.};
\node[font=\tiny\itshape,text=fundgray,anchor=north west,text width=118mm] at (-1.30,-3.55)
  {Thursday is filled because it is the day that makes the following week
  possible; every other day spends what Thursday built. Seven and a half hours a
  week in total.};
\end{appfig}
\vspace{-0.60cm}
\figcaption{Figure 14. The weekly cadence as five standing functions with a time budget\\
each, and the carry rule that moves unfinished work a week rather than a day.}
\end{appfloat}

\subsection{The three rules that make it survive a real week}

A holiday moves the theme rather than deleting it: whichever weekday is closed
becomes the preparation day, and the themes shift around it. That is exactly what
the block just completed did, which is the evidence that the rule works.

No day runs over its budget. A Monday that takes three hours has borrowed from
Thursday, and Thursday is the day that makes the following week possible.

Unfinished work is carried to the same weekday next week rather than squeezed
into the next day. A cadence that reschedules within a week stops being a cadence
by Wednesday.

\subsection{What is deliberately not in the cadence}

Fundraising conversations with no scheduled next action. Reading. Writing the
technical artifacts themselves, which is the rest of the week and is not funding
work \cite{capitalplan2026}. And anything that would be chased rather than waited
for: the four unchaseable items in \S3 are checked on Thursday and never acted
on.

\subsection{How the next block is built from this}

Take the five weekday themes, substitute the current content, and the design work
is done. The block just completed maps onto the frame exactly, in order: federal,
capital, site and partner, preparation, execution. That mapping was not
accidental, and it is the argument that this is the right frame rather than a
convenient one.

\subsection{What the cadence costs, stated honestly}

Seven and a half hours a week is roughly a fifth of a working week spent on
funding rather than on the technical work that produces the artifacts funding
describes. For a company of one person that is a real cost and it is worth
naming rather than presenting the cadence as free.

The alternative is worse in a specific way. Funding work that has no fixed budget
expands to fill whatever anxiety is available, and it does so on exactly the days
when a deadline is near, which are the days the technical work is also most
pressed. A fixed budget converts an unbounded claim on attention into a bounded
one, and the bound is what makes the technical work survivable alongside it.

